\section{近期研究（2026 年 8--11 月）}
\paragraph{以可刷新的语义高斯地图连接移动操作的各阶段}
2026.08~\cite{2608_10756}。该系统将可更新的语义高斯作为主动观测、语言目标定位、障碍几何、可达站位与视觉语言动作模型的共享接口。3DGS 在此承担可查询地图，定位结果还需经几何与可达性处理后才能用于操作。

真实机器人实验考察长流程任务、高度变化、易混淆干扰和杂乱环境。重建、语义定位与低层动作各有误差来源，结果对应论文验证的装置与任务，不能等同于对开放环境的通用操作能力。

\paragraph{SpotlessGS：把车载或机载局部照明与表面反射分开}
2026.08~\cite{2608_14713}。在暗环境中移动自带光源会改变同一表面的外观，普通高斯难以分辨光照变化与场景属性。SpotlessGS 在 DarkGS 基础上给高斯加入法向和可学习反射属性，用轻量 BRDF 网络建模非朗伯行为，同时估计球谐环境光以解释照明。

论文用合成和自采有参考的图像测试重照明，在自采数据上还比较原图与重照明图的 COLMAP 重建。下游结果针对该采集条件，不能等同于已验证长期自主机器人跟踪。

\paragraph{GaussMemory：以任务查询读取并更新持久高斯记忆}
2026.08~\cite{2608_14986}。GaussMemory 把多视角 RGB 经前馈高斯编码转换为观测标记，并在统一注意力中让任务查询读取持久场景记忆。另一输出分支学习把新观测匹配、插入或裁剪记忆原语，使被操作物体的状态能随时间更新。

作者在 LIBERO 的空间、对象、目标和长程任务以及 VLABench 比较成功率和进展。能够处理这类任务依赖多视角感知与训练设置，表中的部分基线还使用额外腕部相机和本体状态，比较时应保留输入差异。

\paragraph{OccamView：在相同采集帧预算下优先补足物体附近几何}
2026.08~\cite{2608_16499}。OccamView 从在线 RGB-D 测量和目标检测维护物体记忆，将附近未观测体素作为物体遮挡区的代理。候选视角按考虑遮挡的代理覆盖率重新排序，但只有原始几何收益满足阈值的视角才允许重排，从而保留原有探索目标。

在 Replica 与 MP3D，作者按相同观测帧数而非规划轮数比较重建，报告完整度收益。遮挡代理不是真实未观测表面，检测遗漏和深度支持不足会影响它的质量。

\paragraph{GS-VLA：在冻结机器人策略前校正相机视点}
2026.08~\cite{2608_19066}。GS-VLA 接在冻结视觉语言动作策略前，先利用标定深度把可见内容几何映射到规范视点，只用局部高斯模块补足新露出的区域。策略权重保持不动，修正发生在图像观测空间。

作者在 LIBERO 四个套件对比多种策略与相机扰动幅度。方法需要单视图的标定深度，在源视图完全看不到的内容不能由重投影可靠复原，因此不能把有界扰动的结果推成大幅换视角鲁棒性。

\paragraph{ChainSplat：以开链动力学连接可变形线体的高斯外观}
2026.08~\cite{2608_28570}。ChainSplat 将线状物的运动表示为基于旋量理论的低维开链，同时让每段链节携带高斯外观。从多视图 RGB 先联合恢复外观、几何与关节状态序列，再识别动力学参数；这个物理模型可用于后续状态估计和模型驱动动作优化。

论文在桌面接触和空中悬挂两类机器人操作中评估，分别对应摩擦和重力主导的条件。开链和多视角观测是方法假设，不能直接推广到拓扑变化或任意软体。

\paragraph{RIDE：用高斯重定位的米制锚点校准单目视频深度}
2026.09~\cite{2609_11079}。RIDE 从既有米制高斯场执行重定位，保留通过 PnP-RANSAC 的匹配，用这些稀疏米制深度约束冻结的视频深度模型。在线分开估计全局尺度、有限空间修正和时序记忆，使少量锚点不必单独决定全部像素深度。

论文按场景划分公共 RGB-D 数据，并测试五个场景中的 27 条真实机器人路线。部署不需额外深度传感器来提供每帧锚点，但前提是已有米制高斯地图且能找到有效重定位匹配。

\paragraph{Neverwhere：从真实场景重建构建视觉跑酷基准}
2026.09~\cite{2609_16443}。Neverwhere 从多视图影像计算相机位姿和多视立体网格，再采样网格初始化高斯。网格承载仿真场景几何，高斯提供视觉观察。构建场景时仍需手工对齐坐标并标注跑酷路线点。

基准把跨障碍、越间隙、斜坡和楼梯拆成逐类任务，检查视觉运动策略成功率和路线进展。这里重建与真实环境之间的差异仍需通过实机试验衡量，基准上的泛化不能直接推出部署表现。

\paragraph{SemSafe-3DGS：用语义和地图不确定性调整导航安全约束}
2026.09~\cite{2609_19330}。SemSafe-3DGS 给地图原语增加语义风险权重，把碰撞后果和几何不确定性一起用于构造控制屏障。主动观测根据当前轨迹调整感知方向，以改善真正影响控制的局部地图。

实验分别考察组合屏障效率、轨迹相关主动感知，以及 Ackermann 实机部署。安全约束依赖地图、语义和所用车辆动力学，观察到的导航收益不等同于任意场景的无碰撞保证。

\paragraph{AirSplan：在高斯场中评估飞行路径的碰撞风险}
2026.09~\cite{2609_21226}。AirSplan 用一组随轨迹变化的体积包络近似飞行器占据，在归一化高斯地图上估计这些包络的碰撞风险，并把风险作为多候选路径优化的约束。这里高斯地图是风险查询的输入。

论文分别在仿真中比较碰撞查询和端到端规划，作者明确指出实机验证仍缺，当前还假定地图准确且能追踪计划轨迹，并需要外部计算。结果不能视作实飞安全保证。

\paragraph{Splat-CBF：在同一控制问题中处理避障与主动观测}
2026.09~\cite{2609_23100}。Splat-CBF 把风险感知的安全屏障与沿轨迹的信息感知屏障放进同一个控制二次规划。发生冲突时，用自适应松弛降低感知要求，并保留避障约束。

四场景的受控模拟中，两种配置观测到的安全率均为 99.5\%，与 SAFER-Splat 相同。论文还区分室内导航、Isaac Sim 机械臂与实物 Ackermann 车辆。所用动力学、地图与软硬约束有别，不能直接交换三组结果。

\paragraph{PlenoCI：用全光场导数区分几何和外观变化}
2026.09~\cite{2609_28930}。两次建图的颜色差可能是物体真的变了，也可能是观察方向改变。PlenoCI 将高斯场用作全光场，用可解析的空间和方向导数刻画原语在不同射线下的响应。实例高斯提供对象组织，导数特征再用于发现变化和区分几何、外观两类原因。

论文将 CL-Splats 和 PASLCD 的二类变化检测与 PASLCD 的多类评价分开，并用未变化独立采集考察误报。对照属于所测静态与半静态场景，不能从几何与外观的区分推断任意动态光照或运动都能正确分类。

\paragraph{Robot-GST：将重建外观和交互几何联合用于机器人行为预演}
2026.09~\cite{2609_33872}。Robot-GST 用短 RGB-D 记录和 SAM3D 对齐可渲染高斯及对象级几何代理。高斯用于提供照片级外观，显式代理承担碰撞、接触等查询。语言任务先变成以三维关键点和关系表示的多段行为，再在重建环境里比较候选结果。实机执行使用估计的目标状态规划动作。

作者用刚性、软质和可形变物体任务检查仿真筛选与实机表现。遮挡、场景变化及接触模型近似会影响预演可靠性，照片级渲染也不等于已经验证相应物理参数。

\paragraph{CollisionSplatting：用高斯空间不确定性评估避碰代价}
2026.09~\cite{2609_35619}。CollisionSplatting 把场景高斯的空间不确定性纳入距离和避碰代价，利用可调保守系数扩大对潜在障碍的避让。GPU 并行碰撞检查作为模块接入采样规划和模型预测控制，视觉渲染特征可同时用来定义目标。

论文给出仿真中的吞吐量、显存和路径对比，也在真实机械臂和移动机器人上展示了受视觉目标约束的任务。保守度提高通常会带来更多绕行，图中以高斯中心的最近距离做代理指标，实验结果不能当作对新场景无碰撞的保证。

\paragraph{从带状态的模拟教师蒸馏仅用 RGB 的导航安全滤波}
2026.09~\cite{2609_36520}。该工作用高斯渲染生成的交互仿真为教师提供动态障碍状态。训练时，控制屏障函数教师修改给定导航指令，学生从相机图像序列与控制输入学习类似校正。部署时不用查询三维高斯地图，地图在这里主要服务训练和模拟。

作者在模拟和实机动态避障试验中统计成功率、无碰撞率与安全距离。学生的表现由这些试验衡量，不能把教师使用的屏障函数形式自动解释成实机学生的形式安全证明。

\paragraph{在返航时间约束下主动采集三维高斯重建视角}
2026.10~\cite{2610_04013}。系统用 RGB-D 更新在线高斯地图供渲染重建，并另保留 TSDF 与占据数据作为碰撞查询和三维 Voronoi 路网的基础。下一视角选择同时考虑信息量、行程和剩余飞行时间，并保留返航所需的时间预算。

作者比较室内照片级仿真中的覆盖、重建误差和返航结果，Gibson 与 HM3D 用场景平均，ReplicaCAD 用物体实例平均。周期性重建路网和遇障时的反应式调整仍有时间开销，返航表现只在所测试的仿真预算与航行模型下得到评价。

\paragraph{FOCUS：从既有高斯地图分辨物体保留、替换和移除}
2026.10~\cite{2610_05639}。同一位置换成外形相似的新物体时，只比对占据体积容易漏检。FOCUS 用带位姿的 RGB-D 序列建立参考图，在无变化的首次采集回放中估计各物体的正常观测残差。后续访问将几何一致性与高斯渲染出的视角相关外观作为两类似然，连续更新保留、替换和移除三种状态。

作者在 Isaac Sim 仓库中报告物体替换 F1 从概率基线的 0.20 升至 0.84，并以相同参数评测真实 TorWIC。消融显示，几何通道能识别移除但难以识别原地替换，单用外观又更易误报。结果依赖参考建图、分割和重访观测，不足以保证换季或长期光照变化下仍能保持原有校准。

\paragraph{GS-Pool：在两次高斯建图中汇总逐对象的变化证据}
2026.10~\cite{2610_06688}。GS-Pool 需要对同一空间两次访问分别训练高斯场，并使相机坐标一致。方法将 SAM2 的二维对象掩码提升到有渲染贡献的高斯，在对象内汇总几何、颜色、语义与光度变化线索。光度线索来自把另一次访问的照片作为对照，将图像误差经渲染器回传到高斯。共享未变对象提供每个场景自己的参照。

在 PASLCD 的共同坐标设置下，采用 FastPGSR 的一组结果为 mIoU 0.7509、F1 0.8461，作者报告 mIoU 相对既有最佳结果提高约 17\%。换用原始 3DGS 训练器时，论文报告相对提升约 12\%，两者是不同配置下的 mIoU 对照。分割未产生单独掩码的变化无法在后续找回，单独光照变化也可能造成误报。

\paragraph{SURGE：融合视觉和声呐定位后进行目标级高斯重建}
2026.10~\cite{2610_07472}。水下单目视觉能提供目标成像方向，但缺少可靠米制尺度。SURGE 利用图像门控从多帧二维成像声呐中获取距离线索，在因子图里联合估计拖缆机器人的轨迹和目标中心，并让声学残差容许主要回波来自近侧表面。随后用米制位姿训练目标范围内的声呐高斯。

实验使用淡水中环绕单个目标的记录，按离线流程处理，报告相对光学基线更集中的目标表示。声呐俯仰歧义、视角覆盖有限和缺少独立测量真值限制了结论，结果并未验证精确表面网格或全场景重建。

\paragraph{2DGS-Planner：直接在可渲染地图中查询路径净空}
2026.10~\cite{2610_11752}。2DGS-Planner 不把整张高斯地图先转换成密集占据数组，而对候选路线调用特定视点的光栅化。多视图结构证据分配路线图节点，候选边附近的几何查询形成可缓存的局部净空数据，再用于不同机器人半径下的在线路径搜索。

原文对路线图构建、净空保真和完整路径规划分开评价。定量环境包括拥有真实参考网格的合成 Stonehenge 与只提供配准代理网格的 Tokyo，两处几何参照可信度不同。系统的净空判断仍取决于重建覆盖与地面估计，实验主要测已有静态地图。

\paragraph{NovaSplat：从手机视频建立无 CAD 对象位姿模型}
2026.10~\cite{2610_13309}。NovaSplat 用目标物体短视频进行一次建模，经 SfM、多视立体和高斯训练获得几何模板和可微渲染表示。查询阶段先做候选定位与外观、几何匹配，再固定高斯参数，通过光度和几何损失联合优化物体相对相机的六自由度位姿。

HOPEv2 和 IndustryShapes 分别按 BOP 的检测和定位协议度量，不能横向比较两者的 AP 与召回。作者报告在 HOPEv2 上与使用 CAD 模型的方法比较有优势，在 IndustryShapes 上仍有差距，尤其是低纹理和强反光物体会影响重建与匹配。

\paragraph{YUBI-Real2Sim：检验高斯外观缺陷会不会改变策略排名}
2026.10~\cite{2610_17521}。YUBI-Real2Sim 用实机排名作为参照，在物理、物体、初值和随机种子保持一致时改动高斯背景，分开比较空洞、漂浮伪影和减少采样的设置。这样观察到的排名变化与外观条件有较清楚的因果联系。

清洁重建的仿真在九个场景任务组合中有八个恢复实机排名，而缺陷条件下更少。排名即使不变，也可能失去统计区分。依据来自所测场景和任务，不能单凭高图像分数保证任意策略排序。

\paragraph{I2O2I-Nav：将室内外连续导航拆成阶段并单独评价交接}
2026.10~\cite{2610_18307}。I2O2I-Nav 把出楼、室外移动、入楼和开放词汇目标搜索串成连续任务。场景用高斯提供视觉，用配准几何判断碰撞。其模块化智能体协调原有子策略，分别处理完成确认、阶段交接和有限次数恢复。

实验分别评估单阶段、一次交接及完整任务，在二十个场景的 1374 个回合用受碰撞约束的运动学智能体运行。人工布景和重建伪影会影响基准覆盖与感知，结果尚不包含关节动力学条件的实机验证。

\paragraph{GS-PRM：用高斯渲染构建和查询视觉路网}
2026.10~\cite{2610_19541}。GS-PRM 以高斯渲染的深度、累计不透明度及覆盖度判断候选位置和路径是否得到地图支持，在通过检查的区域建立可复用路网。查询时用目标图像在已有地图中定位，将起终点接入路网，并为局部策略渲染途中目标。

作者在仿真里用真实碰撞几何检查渲染判据，并在真实差速机器人上评价视觉导航。论文给出的完备性和物理无碰撞结论以覆盖、局部可连接性及保守表面渲染假设为前提，低置信地图区域不可据此视作已验证安全。

\paragraph{借助可微高斯场联合设计编码孔径和深度感知}
2026.10~\cite{2610_28004}。该工作把高斯模型作为可微拍摄环境，在候选相机参数下渲染针孔 RGB 与深度，再用随深度变化的点扩散函数模拟编码孔径。恢复全焦图和估计深度的任务损失可反传给相机和孔径参数，使传感器设计与感知网络一起优化。

论文在单目和双目散焦深度任务比较，并补充原型实拍。结果针对所用光学近似、虚拟场景覆盖与评价任务，不能将图像保真度直接解释为所有下游任务的传感器收益。

\paragraph{Real2Sem2Real：从一个视频获得可迁移的物体级动作数据}
2026.10~\cite{2610_28880}。Real2Sem2Real 用带语义特征的高斯表示跟踪示范视频中的物体状态，再将抓取和放置等操作对齐到同类但外形不同的物体，生成可用于仿真的动作数据。流程还需要示范物体扫描和同类物体参考，单个视频提供的是动作示范。

作者比较四类任务和实机回合，也通过相同视频与同规模样本比较是否引入同类物体变体。流程要求示范物体具有可靠米制信息，这项条件限制了从纯单目视频直接执行完整流程。

\paragraph{AOV-GS：早期用语义新颖度选视，后期回到几何细化}
2026.10~\cite{2610_29374}。AOV-GS 在探索初期用冻结 SAM 与 CLIP 提取候选掩码的语义新颖度，让机器人较早观察更多对象；随后回到几何优先的选视策略。探索结束后再用采集到的嵌入训练场景专属的 LangSplatV2 语言场，因此探索与最终文本查询分成两个阶段。

论文在 Replica 和 Matterport3D 的 Habitat 仿真比较。方法使用相同的真值相机位姿，实验隔离了选视策略的影响，没有同时验证带 SLAM 位姿误差时的效果。平均语义收益也不表示每个场景都改善。

\paragraph{AirPlayground：用重建外观和几何统一比较飞行策略}
2026.10~\cite{2610_29699}。AirPlayground 将每个环境的三维高斯用于相机外观，将同坐标系的米制点云转为用于采样与碰撞的几何。通过统一接口和适配器，不同导航策略可使用相同航迹训练，并在相同闭环条件下执行。

论文把 94 个环境划为 70 个训练、8 个验证及 16 个未见测试场景，并分析离线预测与闭环执行的差异。场景的飞行空间仍由点云、占据和人工标注的范围决定，重建仿真得分不能代替真实飞行测量。

\paragraph{SceneAgent：将拍摄记录转换为可交互仿真场景}
2026.10~\cite{2610_34745}。SceneAgent 调度不同的重建与物理建模工具，将物体、背景、关节和接触等表示组装在仿真器里，再用渲染检查识别并修复部分不一致。对于动态视频，还将模拟运动拟合到观测。高斯等视觉重建给出外观线索，但交互所需的几何与物理属性要另外建模。

论文在静态场景和 4DCodeBench 比较重建，并给出初步实机迁移。视觉检查可以发现部分外观或几何问题，却无法确认未观测结构与物理参数。反复调用模型和外部计算也限制复杂场景的转换速度。

\paragraph{PrecisionReach-GS：在机器人坐标下生成腕部观察与接近动作}
2026.10~\cite{2610_35626}。PrecisionReach-GS 把真实工作区的高斯模型对齐到机器人基座坐标，用可行末端位姿渲染腕部相机观测，并和对应的平滑接近轨迹配成视觉动作样本。部署时由腕部相机驱动微调后的视觉语言动作策略，并用指令相关的验证器判断何时停止。

实验在 Franka FR3 的主板场景验证三个部件的接触前接近。中等程度的位置变化与完整布局变化是不同测试，后者可能要重新重建和生成数据；是否能到达目标也不直接证明接触后拆解或抓取的成功。
